\documentclass[text05.tex]{subfiles}
\begin{document}
\section{赤外線動作ロボット(Sikaye-RB1)に関する\ruby{備考}{び|こう}} \label{ir_robot_notice}
赤外線ロボット(Sikaye-RB1)では、赤外線信号の1周期の長さはすべて 33 で、
信号長の\ruby{制限}{せい|げん}を\ruby{超}{こ}えません。
ただし、リモコンのボタンを\ruby{押}{お}している間、信号を\ruby{繰}{く}り返し送信するコマンドがあります。例えば、前進する Forward はそのようなコマンドです。ボタンを長く\ruby{押}{お}し続けると、記録される信号全体も長くなることがあります。
赤外線ロボットの信号を記録するときは、ボタンを必要以上に長く\ruby{押}{お}さず、短く\ruby{押}{お}して\ruby{離}{はな}すようにしましょう。

参考として、赤外線ロボット (Sikaye-RB1) の信号を mode2 コマンドで\ruby{実際}{じっ|さい}に記録したときの信号長を、表 \ref{ir_robot_signal} に\ruby{示}{しめ}します。
この表の長さは、記録を開始してから\ruby{終了}{しゅう|りょう}するまでに\ruby{含}{ふく}まれる信号全体の長さです。そのため、リモコンのボタンを\ruby{押}{お}し続ける時間が長いほど、信号全体も長くなります。
一方、記録した信号から\ruby{繰}{く}り返しの1周期分だけを取り出すと、その長さは表 \ref{ir_robot_signal_ideal} のように、すべて 33 になります。

\begin{table}[htbp] \begin{center}
    \caption{記録される赤外線ロボット(Sikaye-RB1)の信号長の一例}
    \label{ir_robot_signal}
    \begin{tabular}{cr}
\hline
Command & Length \\
\hline
Backward & 136 \\
Forward & 238 \\
good habit & 34 \\
Machine Language & 34 \\
MEMO & 34 \\
Mode Switch rightpad & 34 \\
Mode Switch & 70 \\
Music & 34 \\
Program & 102 \\
Science Popularization & 34 \\
Slide Backward & 68 \\
Slide Forward & 34 \\
Song & 34 \\
STOP & 34 \\
Turn Left rightpad & 102 \\
Turn Left & 68 \\
Turn Right rightpad & 102 \\
Turn Right & 102 \\
Volume minus & 34 \\
Volume plus & 34 \\
\hline
    \end{tabular}
\end{center} \end{table}

\begin{table}[htbp] \begin{center}
    \caption{赤外線ロボット(Sikaye-RB1)の信号の一周期分の長さ}
    \label{ir_robot_signal_ideal}
    \begin{tabular}{cr}
\hline
Command & Length \\
\hline
Backward & 33 \\
Forward & 33 \\
good habit & 33 \\
Machine Language & 33 \\
MEMO & 33 \\
Mode Switch rightpad & 33 \\
Mode Switch & 33 \\
Music & 33 \\
Program & 33 \\
Science Popularization & 33 \\
Slide Backward & 33 \\
Slide Forward & 33 \\
Song & 33 \\
STOP & 33 \\
Turn Left rightpad & 33 \\
Turn Left & 33 \\
Turn Right rightpad & 33 \\
Turn Right & 33 \\
Volume minus & 33 \\
Volume plus & 33 \\
\hline
    \end{tabular}
\end{center} \end{table}
\end{document}
